\documentclass[11pt]{article}
\usepackage{mathunal-base}
\mathunalfuente{serif}
\usepackage{listings}

\renewcommand{\examtitulo}{Segundo Examen Parcial de Álgebra Lineal - A}
\renewcommand{\examfecha}{29 de julio 2019}

\newcommand{\vf}{\fbox{V}\ \fbox{F}}

\begin{document}

\examheader

\vspace{4pt}
\noindent
\begin{minipage}[t]{0.72\linewidth}
\textbf{Puntaje. Sólo para uso Oficial}\\[4pt]
\begin{tabular}{|c|c|c|c|c|c|}
\hline
I. 1-7 & II. 1 & II. 2 & II. 3 & \textbf{TOTAL} & \textbf{NOTA} \\
(46) & (18) & (20) & (16) & (100pts) & (5.0=100pts) \\ \hline
 & & & & & \\ \hline
\end{tabular}
\end{minipage}

\vspace{6pt}
\noindent\textbf{Instrucciones:} La duración del examen es de 1 hora y 50 minutos. El examen consta de 7 preguntas en la parte de completación y tres preguntas en la parte de procedimiento, en tres hojas impresas por ambos lados. Verifique que su examen esté completo y consérvelo con el gancho. En las preguntas con procedimiento justifique sus respuestas en los espacios asignados. No está permitido sacar hojas en blanco ni ningún tipo de apuntes durante el examen. Verifique que su celular esté apagado y que no esté a la mano. No se permite el uso de calculadora.

\vspace{6pt}
\noindent\textbf{IDENTIFICACIÓN}

\vspace{8pt}
\noindent Nombre: \dotfill\ Cédula \dotfill

\vspace{4pt}
\noindent Profesor: \dotfill\ Grupo \dotfill

\vspace{8pt}
\fbox{\begin{minipage}[t][5cm]{\dimexpr\linewidth-2\fboxsep-2\fboxrule}
\textbf{Espacio para completar respuestas o borrador. Marque claramente si hay respuestas en este espacio.}\\
\textbf{No voltee esta hoja hasta ser AUTORIZADO por el encargado de salón.}
\end{minipage}}

\newpage
\noindent\textbf{\large I. Completación}

\vspace{2pt}
\noindent En las preguntas 1 a 5 complete los espacios en blanco. \textbf{NOTA}: En esta sección se califica sólo la respuesta y no se tiene en cuenta el procedimiento. \textbf{Sólo recibe puntos si la respuesta es completamente correcta.}

\begin{enumerate}
\item{}[8pts] Para cada uno de los siguientes afirmaciones marque con X la letra de su elección.
\begin{enumerate}
\item{}[2pts] La intersección de cualquier plano y de cualquier recta en $\mathbb{R}^3$ es un subespacio. \hfill\vf

\item{}[2pts] La unión de la recta $y=-x$ y $y=5x$ es un subespacio de $\mathbb{R}^2$. \hfill\vf

\item{}[2pts] $\left\{\begin{bmatrix}x\\y\\z\end{bmatrix}\in\mathbb{R}^3: x+z=2x+y+1\right\}$ es un subespacio de $\mathbb{R}^3$. \hfill\vf

\item{}[2pts] $\left\{\begin{bmatrix}x\\y\\z\end{bmatrix}\in\mathbb{R}^3: x=y,\ z=2x+y\right\}$ es un subespacio de $\mathbb{R}^3$. \hfill\vf
\end{enumerate}

\item{}[5pts] Sea $\mathcal{L}$ la recta que pasa por el origen y es perpendicular al plano $x-2y+3z=0$. Calcule una base para el subespacio $\mathcal{L}$.

\vspace{4pt}
\fbox{\begin{minipage}[t][2.8cm]{\dimexpr\linewidth-2\fboxsep-2\fboxrule}
$\mathcal{B} = $
\end{minipage}}

\item{}[8pts] La matriz $A = \begin{bmatrix}1&-1&-1&-1\\-1&1&0&2\\3&-3&-1&3\end{bmatrix}$ tiene forma escalonada reducida $R=\begin{bmatrix}1&-1&0&2\\0&0&1&3\\0&0&0&0\end{bmatrix}$.
\begin{enumerate}
\item{}[3pts] Determine una base $\mathcal{B}_1$ para $\mathrm{Col}(A)$.

\vspace{4pt}
\fbox{\begin{minipage}[t][2.5cm]{\dimexpr\linewidth-2\fboxsep-2\fboxrule}
$\mathcal{B}_1 = $
\end{minipage}}

\item{}[5pts] Determine una base $\mathcal{B}_2$ para $\mathrm{Nul}(A)$.

\vspace{4pt}
\fbox{\begin{minipage}[t][2.5cm]{\dimexpr\linewidth-2\fboxsep-2\fboxrule}
$\mathcal{B}_2 = $
\end{minipage}}
\end{enumerate}

\item{}[5pts] Dadas las transformaciones lineales de $\mathbb{R}^2$ en $\mathbb{R}^2$: $P$ es la proyección sobre el eje $x$, y $R$ es la rotación por un ángulo de 90 grados en el sentido contrario a las manecillas del reloj. Escriba la matriz estándar asociada a la transformación $[R\circ P]$.

\vspace{4pt}
\fbox{\begin{minipage}[t][2.5cm]{\dimexpr\linewidth-2\fboxsep-2\fboxrule}
$[R\circ P] = $
\end{minipage}}

\item{}[5pts] Se quiere extender $\left\{\begin{bmatrix}1&1\\1&1\end{bmatrix},\begin{bmatrix}0&1\\1&0\end{bmatrix}\right\}$ a una base del espacio de matrices simétricas $2\times 2$. Escriba una lista de vectores los cuales agregados a los dados completan una base.

\vspace{4pt}
\fbox{\begin{minipage}[t][2.5cm]{\dimexpr\linewidth-2\fboxsep-2\fboxrule}
\phantom{x}
\end{minipage}}

\item{}[5pts] Sea $B=[I_3\ I_3\ 0_{3,2}]$ donde $I_3$ es la matriz identidad $3\times 3$ y $0_{3,2}$ es la matriz cero $3\times 2$. Determine $\mathrm{Nulidad}(B)$.

\vspace{4pt}
\fbox{\begin{minipage}[t][2.3cm]{\dimexpr\linewidth-2\fboxsep-2\fboxrule}
$\mathrm{Nulidad}(B) = $
\end{minipage}}

\item{}[10pts $=5\times 2$pts] Considere la siguiente secuencia de instrucciones en Matlab:
\begin{quote}
\ttfamily
\noindent >> u = [1; 1; 1; 1]; v = [1; 1; 1; 0]; w = [0; 1; 1; 0]; x = [1 0 1 0]';\\
>> A = [u v w eye(4)];\\
{[}B, p{]} = rref(A);\\
>> C = B(:,1:3);\\
>> D = A(:,p);
\end{quote}
\noindent Escriba, usando nuestra notación matemática usual para matrices, los valores de \texttt{u}, \texttt{x}, \texttt{A}, \texttt{C}, \texttt{D} al final de la ejecución.

\vspace{4pt}
\fbox{\begin{minipage}[t][5cm]{\dimexpr\linewidth-2\fboxsep-2\fboxrule}
\hspace{2cm} $u=$ \hspace{4cm} $x=$

\vspace{2cm}
$A=$ \hspace{3cm} $C=$ \hspace{3cm} $D=$
\end{minipage}}

\noindent \textit{Sugerencia:} Para \texttt{C}, \texttt{D} no necesita realizar la computación de \texttt{rref} completamente, si entiende lo que está pasando.
\end{enumerate}

\newpage
\noindent\textbf{\large II. Solución con Procedimiento}

\begin{enumerate}
\setcounter{enumi}{7}
\item{}[18pts] Sea $\mathcal{B}=\{1,\ x+x^2,\ x^2+x^3,\ x^3\}$ y $\mathcal{H}=gen(\mathcal{B})$.
\begin{enumerate}
\item{}[10pts] Muestre que $\mathcal{H}=\mathcal{P}_3$.

\vspace{6cm}

\item{}[8pts] Para $p(x)=3+5x-2x^2+3x^3$ obtenga su vector de coordenadas $[p(x)]_{\mathcal{B}}$.

\vspace{5cm}
\end{enumerate}
\noindent \textit{Nota:} Las dos partes se pueden resolver en buena parte con el mismo procedimiento, si es el apropiado. En todo caso, asegúrese de marcar claramente su respuesta a cada una de las partes o de lo contrario no recibirá los puntos correspondientes.

\item{}[20pts] Sea $W=\left\{\begin{bmatrix}a&b\\c&d\end{bmatrix}\in M_{2,2}\ \middle|\ \begin{aligned}a+b-c-d&=0\\2a-b+3c+5d&=0\end{aligned}\right\}$.
\begin{enumerate}
\item{}[12pts] Verifique que $W$ es un espacio vectorial.

\vspace{5cm}

\item{}[8pts] Determine una base y la dimensión de $W$.

\vspace{5cm}
\end{enumerate}
\noindent \textit{Nota:} Las dos partes se pueden resolver en buena parte con el mismo procedimiento, si es el apropiado. En todo caso, asegúrese de marcar claramente su respuesta a cada una de las partes o de lo contrario no recibirá los puntos correspondientes.

\newpage
\item{}[16pts] Considere el espacio vectorial $\mathcal{M}_{2\times 2}$, la matriz $C=\begin{bmatrix}1&-1\\-1&1\end{bmatrix}$ y la transformación $T:\mathcal{M}_{2\times 2}\longrightarrow\mathcal{M}_{2\times 2}$ dada por $T(A)=CA-AC$.
\begin{enumerate}
\item{}[10pts] Compruebe que $T$ es una transformación lineal.

\vspace{6cm}

\item{}[6pts] Determine $T$ explícitamente. Es decir, determine $T\begin{bmatrix}a&b\\c&d\end{bmatrix}$ como una matriz $2\times 2$ con cada uno de sus entradas en términos de $a,b,c,d$.

\vspace{5cm}
\end{enumerate}
\end{enumerate}

\examfooter
\end{document}
